%=============================================================================
%  A_appendices.tex  -  Appendices
%=============================================================================
\chapter{Appendices}
\label{ch:appendix}

\section{Environment configuration template}
\label{app:env}
The complete configuration template shipped as \code{.env.example}. Copy it to
\code{.env} and adjust for the target runtime.

\begin{lstlisting}[style=bash,caption={\code{.env.example} (annotated).}]
# --- Server ---
HOST=0.0.0.0
PORT=8000
ENVIRONMENT=development           # development | production

# --- Ollama (LLM runtime) ---
OLLAMA_HOST=http://localhost:11434
# Docker (Ollama on host, backend in container):
# OLLAMA_HOST=http://host.docker.internal:11434

# --- Models ---
MODEL_NAME=gemma4:12b             # primary answering model (text + image)
ROUTER_MODEL_NAME=gemma4:e4b      # fast routing / classification
EMBEDDING_MODEL=all-MiniLM-L6-v2  # RAG embeddings

# --- RAG ---
MAX_CONTEXT_TOKENS=6000           # assembled RAG context budget

# --- Vector store (ChromaDB) ---
CHROMA_PERSIST_DIRECTORY=./chroma_data

# --- CORS ---
ALLOWED_ORIGINS=http://localhost:5173,http://localhost

# --- Multimodal chat ---
IMAGE_CHAT_ENABLED=true
MAX_CHAT_IMAGES=3
MAX_CHAT_IMAGE_MB=5
CONTEXT_WINDOW_TOKENS=256000      # front-end token-meter window
\end{lstlisting}

\section{The chat event protocol}
\label{app:events}
The \code{/chat} endpoint streams Server-Sent Events. Each event has a name and a
JSON \code{data} payload. The front end's \code{useChat} hook switches on the
event name to update the interface incrementally.

\begin{longtable}{@{}L{2.3cm} L{4.6cm} L{7.4cm}@{}}
\caption{Chat SSE events, payloads, and UI effect.}\label{tab:sseevents}\\
\toprule
\textbf{Event} & \textbf{Payload (illustrative)} & \textbf{UI effect} \\
\midrule
\endfirsthead
\multicolumn{3}{c}{\small\itshape (continued from previous page)}\\
\toprule
\textbf{Event} & \textbf{Payload} & \textbf{UI effect} \\
\midrule
\endhead
\midrule
\multicolumn{3}{r}{\small\itshape (continued on next page)}\\
\endfoot
\bottomrule
\endlastfoot
\code{skill} & \code{\{skill, reason\}} & Shows the selected skill. \\
\code{thinking} & \code{\{text\}} & Renders a reasoning affordance. \\
\code{action} & \code{\{action, detail\}} & Indicates a backend action. \\
\code{token} & \code{\{token\}} & Appends to the live answer. \\
\code{chart} & \code{\{spec / figure\}} & Adds a validated dashboard tile. \\
\code{usage} & \code{\{prompt, completion, total\}} & Updates the token meter. \\
\code{command} & \code{\{command, args\}} & Switches view / triggers download. \\
\code{error} & \code{\{message\}} & Displays the real error verbatim. \\
\code{done} & \code{\{\}} & Marks the turn complete. \\
\end{longtable}

\section{The chart-type catalogue}
\label{app:charts}
The twenty-six members of the \code{ChartType} enumeration, the single source of
truth for both the plan schema and the factory's validation.

\begin{lstlisting}[style=py,caption={\code{ChartType} enumeration.}]
BAR, LINE, SCATTER, PIE, HISTOGRAM, HEATMAP, BOX, VIOLIN, STRIP,
FUNNEL, FUNNEL_AREA, TREEMAP, SUNBURST, ICICLE, WATERFALL, AREA,
RADAR, DENSITY_HEATMAP, DENSITY_CONTOUR, SCATTER_GEO, CHOROPLETH,
SCATTER_MATRIX, PARALLEL_COORDINATES, PARALLEL_CATEGORIES,
SCATTER_3D, GAUGE, TIMELINE
\end{lstlisting}

\section{Repository structure}
\label{app:repo}
\begin{lstlisting}[style=bash,caption={Project layout (abridged).}]
DEPI_Graduation_Project_LLM_Driven/
  backend/
    app/
      api/            # chat, upload, sessions, clean, predict,
                      # visualizations, skills, health
      core/           # ingestion, chunker, profiler, session_manager
      eda/            # auto_eda
      llm/            # ollama_client, prompt_templates, token_counter
      rag/            # store, embedder, retriever, context_builder
      skills/         # cleaning, prediction, report, report_narrative
      skills_registry/# markdown skill definitions
      visualization/  # chart_factory, chart_types, export
      models/         # schemas, enums
      config.py, main.py
    tests/            # test_ingestion, test_profiler, test_visualization
    Dockerfile, requirements.txt, pyproject.toml
  frontend/
    src/
      components/     # chat, dashboard, upload, charts, sidebar, palette
      hooks/          # useChat, useSessions, useFileUpload
      lib/            # api, chartExport, reportDownload, tokenUsage
    Dockerfile, vite.config.js, tailwind.config.js, nginx.conf
  scripts/            # setup_ollama.sh, serve_frontend.sh
  docs/               # api_reference.md, deployment.md, latex/
  data/samples/       # example datasets
  main.ipynb          # cloud setup notebook
  docker-compose.yml, docker-compose.prod.yml, Makefile
  .env.example
\end{lstlisting}

\section{Technology summary}
\label{app:tech}
\begin{table}[h]
\centering
\caption{Complete technology stack.}
\small
\begin{tabularx}{\linewidth}{@{}L{3.2cm} X@{}}
\toprule
\textbf{Layer} & \textbf{Technology} \\
\midrule
Primary LLM & Gemma\,4 12B via Ollama (multimodal). \\
Router model & Gemma\,4 \code{e4b}. \\
Embeddings & \code{sentence-transformers} (\code{all-MiniLM-L6-v2}). \\
Vector store & ChromaDB (persistent, per-session collections). \\
Backend & FastAPI, Python\,3.12, \code{pydantic-settings}, \code{uvicorn},
\code{httpx}. \\
Data & pandas, numpy, pdfplumber. \\
ML & scikit-learn. \\
Visualization & Plotly, Kaleido. \\
Reports & ReportLab. \\
Frontend & React\,18, Vite, TailwindCSS, react-plotly.js, react-markdown,
lucide-react. \\
Infra & Docker Compose (dev + prod/nginx), Makefile, GitHub Actions, Ruff,
pytest. \\
\bottomrule
\end{tabularx}
\end{table}

\section{Glossary}
\label{app:glossary}
\begin{description}[style=nextline,leftmargin=0pt,itemsep=3pt]
  \item[\textcolor{uaIndigoD}{RAG}] Retrieval-Augmented Generation --- grounding a
    model's answer in facts retrieved from an external index rather than its
    weights.
  \item[\textcolor{uaIndigoD}{Semantic type}] The profiler's inferred role of a
    column (boolean, datetime, id, continuous, categorical, text), beyond its raw
    dtype.
  \item[\textcolor{uaIndigoD}{Chart plan}] A schema-constrained JSON description
    of a chart emitted by the model and validated before rendering.
  \item[\textcolor{uaIndigoD}{Skill}] A routed capability (analysis,
    visualization, cleaning, prediction, report, Q\&A) with Markdown guidance and
    a Python implementation.
  \item[\textcolor{uaIndigoD}{No-fabrication rule}] The commitment that the system
    shows real model output or a real error, never a placeholder.
  \item[\textcolor{uaIndigoD}{Structured output}] Ollama's mode for forcing
    JSON- or schema-valid model responses.
  \item[\textcolor{uaIndigoD}{Session}] One dataset-scoped analytical context,
    holding the DataFrame, profile, chat history, and generated charts.
\end{description}

\vfill
\begin{center}
\begin{tikzpicture}
  \node[font=\itshape\small,text=uaSlate]
    {Universal Analyst \;\textcolor{uaAmber}{\textbullet}\; Technical Documentation \;\textcolor{uaAmber}{\textbullet}\; DEPI Graduation Project};
\end{tikzpicture}
\end{center}
